% from ega4-11-fr.tex
\oldpage[IV-3]{156}

Jacobson~; on se réduit ainsi au cas où $k'$ est une extension \emph{finie} de
$k$, ce que l'on va supposer dans ce qui suit. Soient $k''$ la plus grande
extension séparable de $k$ contenue dans $k'$, $k_1$ une extension galoisienne
finie de $k$ contenant $k''$, de sorte que $k''\otimes_k k_1$ est composée
directe de corps isomorphes à $k_1$~; comme $k'$ est une extension radicielle de
$k''$, $k'\otimes_k k_1$ est donc composée directe d'extensions radicielles de
$k_1$. Il existe un anneau local $A_1$ qui est une $A$-algèbre et un $A$-module
\emph{libre de type fini}, tel que $A_1/\mathfrak m A_1$ soit $k$-isomorphe à
$k_1$ $(\mathbf 0_{\mathrm{III}},10.3.1.2)$~; de façon précise, on peut supposer
que $k_1=k[T]/(r)$, où $r$ est un polynôme unitaire irréductible et séparable de
$k[T]$ de degré $n$~; si $R$ est un polynôme unitaire de $A[T]$ dont l'image
canonique est $r$ (et qui est donc de degré $n$), on peut prendre
$A_1=A[T]/(R)$. Or, si $K$ est le corps des fractions de $A$, il est clair que
$R$ est un polynôme irréductible et séparable de $K[T]$~; on en déduit donc
d'abord que $A_1$ est un anneau intègre dont le corps des fractions
$K_1=K[T]/(R)$ est une extension séparable de $K$. En outre, si $t$ est l'image
canonique de $T$ dans $A_1$, les $t^j$ $(0\leq j<n)$ forment une base du
$A$-module $A_1$, et leurs images dans $k_1$ une base sur $k$~; on déduit de là
que
\[
d=\det\bigl(\operatorname{Tr}_{A_1/A}(t^{i+j})\bigr)
\]
est un élément de $A$ dont la classe dans $k$ est $\ne0$, et qui est par suite
inversible. Le même raisonnement que dans (6.12.4.1, I)) prouve alors que le
morphisme $\operatorname{Spec}(A_1)\to\operatorname{Spec}(A)$ est plat et a ses
fibres régulières~; on conclut par suite de (6.5.4, (ii)) que $A_1$ est
\emph{intégralement clos}. Considérons alors l'anneau
$A'_1=A'\otimes_A A_1$~; c'est un $A'$-module libre de type fini, donc un anneau
semi-local (Bourbaki, \emph{Alg. comm.}, chap.~IV, \S~2, n$^\circ$~5, cor.~3 de
la prop.~9)~; en outre, les idéaux maximaux de cette $A'$-algèbre finie sont
tous au-dessus de l'idéal maximal $\mathfrak m'$ de $A'$, et \emph{a fortiori}
contiennent $\mathfrak m A'_1$. Mais
\[
A'_1/\mathfrak m A'_1=(A'/\mathfrak m A')\otimes_k k_1,
\]
et comme $k_1$ est une extension séparable et finie de $k$, le radical de
$A'_1/\mathfrak m A'_1$ est égal à
$(\mathfrak m'/\mathfrak m A')\otimes_k k_1$ (Bourbaki, \emph{Alg.},
chap.~VIII, \S~7, n$^\circ$~2, cor.~2 de la prop.~3)~; si $\mathfrak n_i$
$(1\leq i\leq r)$ sont les idéaux maximaux de $A'_1$, les corps
$A'_1/\mathfrak n_i$ sont donc les corps composant l'algèbre
$k'\otimes_k k_1$, autrement dit ce sont des extensions finies
\emph{radicielles} de $k_1$. D'ailleurs, comme $A\to A'$ est un homomorphisme
injectif, il en est de même de $A_1\to A'_1$, $A_1$ étant un $A$-module plat~;
l'homomorphisme canonique
\[
A'_1\longrightarrow\prod_{i=1}^{r}(A'_1)_{\mathfrak n_i}
\]
étant lui aussi injectif (Bourbaki, \emph{Alg. comm.}, chap.~II, \S~3,
n$^\circ$~3, th.~1), il en est de même du composé de $A_1$ dans ce produit.
Mais $A_1$ est intègre, et les noyaux des homomorphismes
$A_1\to(A'_1)_{\mathfrak n_i}$ sont en nombre fini~; comme leur intersection
est nulle, c'est que l'un d'eux déjà est nul. En d'autres termes, il y a un
$B_1=(A'_1)_{\mathfrak n_i}$ tel que l'homomorphisme $A_1\to B_1$ soit
\emph{injectif} et \emph{local}. Posons $Y'_1=\operatorname{Spec}(B_1)$,
$X'_1=X'\times_{Y'}Y'_1$~; $\mathcal F'_1=\mathcal F'\otimes_{Y'}Y'_1$ est
$Y'_1$-plat en tous les points de $X'_1$ au-dessus de $x'$~; d'ailleurs l'idéal
maximal de $B_1$ est le seul au-dessus de $\mathfrak m'$, par suite tous ces
points ont pour projection dans $Y'_1$ le point fermé $y'_1$. Soit $x'_1$ un de
ces points. Posons d'autre part $Y_1=\operatorname{Spec}(A_1)$,
$X_1=X\times_Y Y_1$, $\mathcal F_1=\mathcal F\otimes_Y Y_1$~; si $x_1$ est la
projection de $x'_1$ dans $X_1$, la projection de $x_1$ dans $X$ est $x$ et sa
projection dans $Y_1$ est le point fermé $y_1$. Si l'on prouve que
$(\mathcal F_1)_{x_1}$ est un $\mathcal O_{y_1}$-module plat, il en résultera
que $\mathcal F_x$ est un $\mathcal O_y$-module plat~; en effet
$\mathcal O_{y_1}$ est un $\mathcal O_y$-module plat, donc
$(\mathcal F_1)_{x_1}$ est un $\mathcal O_y$-module plat
$(\mathbf 0_{\mathrm I},6.2.1)$.
% IV3_B23_P156_END

